% methods.tex -- Experimental methods (paragraph-level headings under Materials and Methods)
For each experiment, we followed the stochastic attention pipeline
of~\cite{varnerSAProtein2026} without modification to the core sampler. Seed alignments were
downloaded from InterPro/Pfam in Stockholm format. Alignments were cleaned by removing columns
with $>50\%$ gaps and sequences with $>30\%$ gaps in the remaining columns. The cleaned alignment
was one-hot encoded ($20$ amino acid channels per position, gaps mapped to zero vectors) yielding
a $20L \times K$ matrix, reduced by PCA (retaining $95\%$ of variance) to a $d \times K$ matrix,
and column-normalized to unit norm. ULA chains were run with step size $\alpha = 0.01$ for
$T = 5{,}000$ steps, initialized near random stored patterns with small Gaussian perturbation.
Samples were collected after a burn-in of 2{,}000 steps with thinning every 100 steps. For
multiplicity-weighted runs, $\beta^{*}(\vr)$ was found by sweeping 50--60 logarithmically spaced
$\beta$ values in $[0.1,\, 500]$ and locating the inflection point of the weighted attention
entropy.

\paragraph{Protein families and functional splits.}
We tested three Pfam families: Kunitz domains (PF00014, $K=99$, $L=53$), SH3 domains (PF00018,
$K=55$, $L=48$), and WW domains (PF00397, $K=420$, $L=31$). For each family, we defined a binary
functional split based on a single marker position. For Kunitz, the P1 position---the primary contact residue with serine proteases, known from
structural and biochemical studies to determine inhibitory specificity---was identified as the
column with the highest Lys frequency; sequences with K or R at P1 were designated as the
functional subset ($K_{\mathrm{des}} = 32$) based on their known association with trypsin
inhibition. For SH3, the conserved Trp in the hydrophobic binding groove was identified as the
column with intermediate Trp frequency ($0.15 < f_W < 0.85$); sequences with W formed the
designated subset ($K_{\mathrm{des}} = 33$). For WW, the specificity-determining position was the
highest-entropy column in the middle third of the alignment; sequences with the most common
residue formed the designated subset ($K_{\mathrm{des}} = 69$). In each case, the designation
reflects external biological knowledge, not a property discovered by the method.

\paragraph{Multiplicity-weighted sampling.}
For each $\rho$ value, the multiplicity vector was set to $r_k = \rho$ for designated patterns and
$r_k = 1$ for background patterns. Weighted samples were generated using the modified Langevin update
(Eq.~\ref{eq:weighted-ula}). Twenty to thirty independent chains were run per condition, yielding
620--930 sequences after burn-in and thinning.

\paragraph{Evaluation metrics.}
Phenotype fidelity was measured as the fraction of generated sequences displaying the designated
subset's marker residue at the marker position. Attention diagnostics recorded the mean softmax
weight on designated patterns during sampling. The soft P1 score was computed by reconstructing the PCA-space
sample to one-hot space and applying a temperature-scaled softmax ($T=0.2$) to the amino acid
channels at the marker position. Sequence diversity was computed as one minus the mean pairwise
sequence identity among generated sequences (estimated from 300--500 random pairs). Amino acid
composition fidelity was measured by KL divergence between generated and reference amino acid
frequencies.

\paragraph{Fisher separation index.}
For each functional split, pairwise cosine similarities were computed between all unit-norm PCA
vectors within Group~A, within Group~B, and between groups. The pooled within-group mean
$\bar{c}_{\mathrm{within}}$ and standard deviation $\sigma_{\mathrm{within}}$ were computed by
weighting by group size; the separation index $S$ was then computed as defined in
Section~\ref{sec:gap}.

\paragraph{Software.}
All experiments were implemented in Julia~1.12 using NNlib.jl (softmax),
MultivariateStats.jl (PCA), and Distributions.jl. Source code and experiment scripts are available
at \url{https://github.com/varnerlab/SA-Binding-Generation-Study}.
